\chapterrule
\section{CNN：怎样用局部连接与参数共享处理图片？}
\label{intro:19-cnn}

文本有顺序，图片则同时具有宽度、高度和局部空间关系。
CNN 利用这种结构，让同一组特征检测规则在图像不同位置重复工作。

\subsection{直接把图片铺成一列，会出现什么问题？}
\phantomsection\label{partq:vision}

一张 $30\times30$ 的灰度图片有 $900$ 个像素；每个像素是一个数字，当然可以把图片铺成长度 $900$ 的\textbf{行向量}，再送入全连接层：
\[
\vect h^\top=\relu(\vect x^\top W+\vect b^\top).
\]
如果隐藏层有 $1000$ 个输出，那么
\[
W\in\R^{900\times1000}
\]
就需要\textbf{900,000 个权重}，还不包括偏置。更糟的是，单纯铺平会使图片中“相邻像素”的结构不再显式体现在模型里；为了在不同位置识别相同的边缘或图案，全连接层往往要重复学习类似模式。

\subsection{小区域扫描：卷积核到底做了什么？}
\phantomsection\label{concept:kernel}
\phantomsection\label{concept:stride}
\phantomsection\label{concept:padding}
\phantomsection\label{concept:feature-map}

考虑一张 $4\times4$ 的极简图片：
\[
X=\begin{bmatrix}
1&1&0&0\\
1&1&0&0\\
0&0&1&1\\
0&0&1&1
\end{bmatrix}.
\]
我们设计一个 $2\times2$ 的\textbf{卷积核（kernel）}
\[
F=\begin{bmatrix}1&0\\0&-1\end{bmatrix}.
\]
把它覆盖在图片左上角的 $2\times2$ 小区域，对应元素相乘，再全部相加：
\[
1\times1+1\times0+1\times0+1\times(-1)=0.
\]
若向右移动一格，看到的区域是
\[
\begin{bmatrix}1&0\\1&0\end{bmatrix},
\]
运算结果是 $1$。这个核比较窗口左上角与右下角：前一个位置较亮，响应为正；
两者相同，响应为零。这里的 $1$ 是局部亮度差的响应值，不是分类概率。
把各位置的响应排起来，就得到一张\textbf{特征图}（feature map）。

\begin{center}
\tagstructbegin{tag=Figure,alttext={二乘二卷积核在四乘四图像上滑动产生三乘三特征图}}
\begin{tikzpicture}[>=Latex,x=9mm,y=9mm]
\foreach \rr/\cc/\val in {
0/0/1,0/1/1,0/2/0,0/3/0,
1/0/1,1/1/1,1/2/0,1/3/0,
2/0/0,2/1/0,2/2/1,2/3/1,
3/0/0,3/1/0,3/2/1,3/3/1}{
 \filldraw[fill=paper,draw=linegray] (\cc,{3-\rr}) rectangle ++(1,1);
 \node[font=\small] at ({\cc+.5},{3.5-\rr}) {$\val$};
}
\draw[very thick,orange,rounded corners=1pt] (0,2) rectangle (2,4);
\draw[very thick,teal,dashed,rounded corners=1pt] (1,2) rectangle (3,4);
\node[font=\footnotesize,anchor=west] at (0,-.52) {橙框：首个窗口；青色虚框：右移一格};
\draw[->,blue,very thick] (4.25,2.1)--(5.22,2.1);
\node[align=center,font=\small] at (6.9,2.1)
 {$F=\begin{bmatrix}1&0\\0&-1\end{bmatrix}$\\逐格乘加};
\draw[->,blue,very thick] (8.3,2.1)--(9.1,2.1);
\node[align=center,font=\small] at (11,2.15)
 {$\begin{bmatrix}0&1&0\\1&0&-1\\0&-1&0\end{bmatrix}$\\$3\times3$ 特征图};
\end{tikzpicture}
\tagstructend

{\footnotesize 首格结果为 $1-1=0$；右移一格后为 $1-0=1$。
其余位置使用同一 $2\times2$ 核，步幅 1、不补零。}
\end{center}
\begin{center}\captionof{figure}{一个 $2\times2$ 卷积核在 $4\times4$ 图片上滑动，生成 $3\times3$ 特征图}\label{fig:cnn-window}\end{center}


当步幅（stride）为 $1$、没有补零（padding）时，沿每个方向可停放的位置数是 $4-2+1=3$，所以输出是 $3\times3$。上图的九个数字都由同一组参数反复扫描得到。

\begin{warningbox}
\phantomsection\label{concept:cross-correlation}PyTorch 等框架的卷积层实现的是\textbf{互相关}（cross-correlation）：卷积核不翻转就直接做逐项乘积；很多教程也把它称为“卷积”。严格数学卷积会翻转核。两者在训练可学习卷积核时都能表征局部模式。
\end{warningbox}

\subsection{卷积为什么能减少参数？}

$2\times2$ 的核只有\textbf{4 个可学习数字}，但可以在所有位置\textbf{共享}这 4 个数字；不必为每个位置另写一套参数。这叫\textbf{权重共享}。局部感受野让层先看到附近的像素，堆叠多个卷积层后可以逐步组合更大范围的模式。

\subsection{彩色图有三个通道，一个输出怎样算？}
\label{sec:cnn-channel-sum}

RGB 彩色图把每个像素分成红、绿、蓝三个数，也就是三张输入数表。
一个输出位置需要同时看三个通道，所以\textbf{一个输出卷积核包含三个权重切片}，
每个切片与对应通道的窗口逐项乘加，再把三份结果相加。

另取一个 RGB 小窗口，三个通道分别为
\[
X_R=\begin{bmatrix}1&0\\0&0\end{bmatrix},\quad
X_G=\begin{bmatrix}1&1\\0&0\end{bmatrix},\quad
X_B=\begin{bmatrix}0&0\\0&1\end{bmatrix}.
\]
为方便手算，先让三个权重切片都等于前面的 $F$，偏置为零。
三个通道的乘加分别得到 $1,1,-1$，相加得到这个输出位置的值：
\[
y=1+1+(-1)=\boxed{1}.
\]
三个切片只是这次取相同初值，训练时是独立权重。
这一整套核有 $2\times2\times3=12$ 个权重，扫描后产生\textbf{一个输出通道}。
若想产生两个输出通道，就使用两整套核，共 $24$ 个权重；每套都读取全部三个输入通道。

由此推广：输入通道数为 $C_{\rm in}$、输出通道数为 $C_{\rm out}$，
方形核宽为 $k$，则权重数为 $k^2C_{\rm in}C_{\rm out}$。
若每个输出通道还有一个偏置，总参数量为
\[
\boxed{k^2C_{\rm in}C_{\rm out}+C_{\rm out}}.
\]
这个数量由核与通道决定，\textbf{不依赖图片的宽高}；图片变大，会增加扫描位置和计算量。

\subsection{池化又是干什么的？}
\phantomsection\label{concept:pooling}

得到响应图后，还想缩小它，同时保留每个区域内较强的信号，怎么办？
可以把一个区域汇成一个数。例如对 $2\times2$ 小块
\[
\begin{bmatrix}1&7\\3&2\end{bmatrix}
\]
做\textbf{最大池化}（max pooling），只保留最大值 $7$。池化可以缩小特征图并对一定范围内的微小位移增加稳定性；但它也会丢失空间信息，因此现代 CNN 并不一定每层都使用池化。

\par\Needspace{62mm}
\tagstructbegin{tag=Figure,alttext={四乘四特征图按不重叠二乘二窗口做最大池化。左上窗口一七三二取七，右上二零五一取五，左下四零一三取四，右下六二二八取八。输出二乘二矩阵第一行为七五、第二行为四八，位置数量从十六缩为四。}}
\begin{center}
\begin{tikzpicture}[>=Latex,x=8mm,y=8mm]
\foreach \x/\y/\c in {0/2/paper,2/2/linegray!15,0/0/linegray!25,2/0/lavender}{
 \fill[\c] (\x,\y) rectangle ++(2,2);
}
\draw[step=1,draw=linegray] (0,0) grid (4,4);
\foreach \r/\c/\v in {0/0/1,0/1/7,0/2/2,0/3/0,1/0/3,1/1/2,1/2/5,1/3/1,2/0/4,2/1/0,2/2/6,2/3/2,3/0/1,3/1/3,3/2/2,3/3/8}{
 \node[font=\small] at (\c+.5,3.5-\r) {$\v$};
}
\draw[thick,ink] (0,2)--(4,2) (2,0)--(2,4);
\draw[->,very thick,flowcolor] (4.6,2)--(8.0,2) node[midway,above,font=\small] {每块取最大值};
\foreach \x/\y/\v/\c in {9/2/7/paper,10/2/5/linegray!15,9/1/4/linegray!25,10/1/8/lavender}{
 \filldraw[fill=\c,draw=linegray] (\x,\y) rectangle ++(1,1);
 \node[font=\small\bfseries] at (\x+.5,\y+.5) {$\v$};
}
\node[font=\small] at (2,-.7) {$4\times4$ 特征图};
\node[font=\small] at (10,-.7) {$2\times2$ 特征图};
\node[font=\small,text=muted] at (6,-1.6) {窗口 $2\times2$，步幅 2：保留每块最大响应，舍去块内的精确位置};
\end{tikzpicture}
\captionof{figure}{最大池化：四个局部窗口各留下一个数，空间位置从 16 个缩为 4 个}
\label{fig:cnn-max-pooling}
\end{center}
\tagstructend


\subsection{一个典型图像分类模型}

\[
\text{图像}\ \longrightarrow\
\text{卷积+激活}\ \longrightarrow\
\text{下采样}\ \longrightarrow\
\text{更多卷积}\ \longrightarrow\
\text{分类输出}.
\]
CNN 的关键是局部连接与权重共享：同一组参数可以检测不同位置的相似图案。

\chapterchecktasks{19-cnn}{理解与复算}{
\item \textbf{比较：}同一卷积核在多个空间位置共享参数，与把整张图片展平成向量后使用全连接层相比，节省了什么？\chapterinlinehint{比较空间位置增多时，独立权重组数如何变化。}
\item \textbf{手算：}$3\times3$ 卷积核、输入 $3$ 通道、输出 $8$ 通道，不计偏置共有多少权重？\chapterinlinehint{按核面积、输入通道、输出通道三项相乘。}
\item \textbf{预测：}保持核宽、通道数和步长不变，输入图像宽高各增大一倍，卷积参数量与输出位置数分别怎样变化？\chapterinlinehint{分别追踪权重形状与输出网格面积。}
}


\par\smallskip
图片有空间结构，直接展平会浪费这种结构。CNN 利用小核扫描、跨位置权重共享和分层特征提取，让模型更有效地识别局部模式。下一章换一条路线：把一张 $4\times4$ 小图切成 Patch，使图像也能输入 Transformer；为了看清切块顺序，那里会换一张每块内容都不同的小图。
